% Secao 4: Impulso na Carreira Dev, Inteligencia Artificial (LLMs) e Pesquisa

\begin{frame}{Alavancando a Carreira de Desenvolvedor}{\faRocket\ Da Lógica Básica à Engenharia de Software de Alto Desempenho}
  \begin{columns}[t]
    \begin{column}{0.32\textwidth}
      \begin{bluecard}{\faCode\ Full-Stack \& Mobile}
        \tiny
        \textbf{O Coração do Mercado}
        \begin{itemize}\setlength{\itemsep}{0.15em}
          \item \textbf{Web I e II:} APIs RESTful, front-ends reativos e back-ends escaláveis.
          \item \textbf{Mobile:} Apps para smartphones Android e ecossistemas móveis.
          \item \textbf{SQL + NoSQL:} Persistência relacional e bancos de alta escala.
        \end{itemize}
      \end{bluecard}
    \end{column}
    
    \begin{column}{0.32\textwidth}
      \begin{purplecard}{\faDraftingCompass\ Arquitetura \& Testes}
        \tiny
        \textbf{O Diferencial Sênior}
        \begin{itemize}\setlength{\itemsep}{0.15em}
          \item \textbf{Padrões de Projeto:} Código limpo e desacoplado com padrões GoF.
          \item \textbf{Testes e Qualidade:} Automação de testes unitários, integração e TDD.
          \item \textbf{Eng. de Software:} Requisitos e arquitetura sustentável.
        \end{itemize}
      \end{purplecard}
    \end{column}
    
    \begin{column}{0.32\textwidth}
      \begin{greencard}{\faCloud\ DevOps \& Nuvem}
        \tiny
        \textbf{Entrega Contínua}
        \begin{itemize}\setlength{\itemsep}{0.15em}
          \item \textbf{DevOps:} Pipelines de integração e entrega contínua (CI/CD).
          \item \textbf{Docker \& Nuvem:} Implantação e orquestração ágil de contêineres.
          \item \textbf{Segurança:} Práticas de DevSecOps e proteção de dados.
        \end{itemize}
      \end{greencard}
    \end{column}
  \end{columns}
  
  \vspace{0.15cm}
  \begin{statcard}
    \scriptsize \textcolor{TechDark}{\faAward\ \textbf{Vantagem Competitiva:} Formação completa que habilita o egresso a disputar vagas júnior/pleno em qualquer empresa do Brasil ou do exterior.}
  \end{statcard}
\end{frame}

\begin{frame}{A Nova Fronteira: IA, LLMs e Pesquisa Científica}{\faBrain\ Aliando Modelos de Linguagem e Inovação Acadêmica}
  \begin{columns}[t]
    \begin{column}{0.48\textwidth}
      \begin{purplecard}{\faRobot\ Inteligência Artificial \& LLMs}
        \scriptsize
        Na disciplina de \textbf{Inteligência Computacional (80h)} e optativas de IA e Dados:
        \begin{itemize}\setlength{\itemsep}{0.15em}
          \item \textbf{Fundamentos Sólidos:} Machine learning, redes neurais e representação de conhecimento vetorial.
          \item \textbf{Copilotos \& Agentes:} Engenharia de prompt e integração de APIs de LLMs para automação inteligente.
          \item \textbf{Arquiteturas RAG:} Criação de sistemas de busca semântica unindo bancos de dados e LLMs.
          \item \textbf{Auditoria Crítica:} Combate a alucinações com rigor de engenharia de software e testes.
        \end{itemize}
      \end{purplecard}
    \end{column}
    
    \begin{column}{0.48\textwidth}
      \begin{greencard}{\faMicroscope\ Crescimento na Pesquisa Acadêmica}
        \scriptsize
        Portas abertas para Mestrado e Inovação:
        \begin{itemize}\setlength{\itemsep}{0.15em}
          \item \textbf{Bolsas de Pesquisa (PIBIC / PIBITI):} Projetos remunerados orientados por mestres e doutores.
          \item \textbf{IA para Demandas Locais:} Aplicações no agronegócio, recursos hídricos, saúde e educação no semiárido.
          \item \textbf{Publicações Científicas:} Apresentação de artigos em congressos da SBC (ERBASE, CBIE, IHC).
          \item \textbf{Ponte para Pós-Graduação:} Vantagem curricular decisiva para ingresso direto em Mestrados e Doutorados.
        \end{itemize}
      \end{greencard}
    \end{column}
  \end{columns}
\end{frame}
